\documentclass[11pt,a4paper]{article}

\usepackage[utf8]{inputenc}
\usepackage[T1]{fontenc}
\usepackage{lmodern}
\usepackage[margin=2.5cm]{geometry}
\usepackage{amsmath, amssymb, amsthm}
\usepackage{graphicx}
\usepackage{booktabs}
\usepackage{siunitx}
\usepackage{listings}
\usepackage{xcolor}
\usepackage{fancyhdr}
\usepackage{enumitem}
\usepackage[hidelinks]{hyperref}

% ----- Course details -----
\newcommand{\coursecode}{TCS1234}
\newcommand{\coursename}{Course Name}
\newcommand{\assignmenttitle}{Assignment 1}
\newcommand{\studentname}{Your Name}
\newcommand{\studentid}{1234567890}
\newcommand{\lecturer}{Dr. Lecturer Name}

\pagestyle{fancy}
\fancyhf{}
\lhead{\coursecode: \assignmenttitle}
\rhead{\studentname}
\cfoot{\thepage}

\theoremstyle{definition}
\newtheorem{question}{Question}
\newenvironment{solution}{\begin{proof}[Solution]}{\end{proof}}

\lstset{
  basicstyle=\ttfamily\small,
  keywordstyle=\color{teal}\bfseries,
  commentstyle=\color{gray}\itshape,
  stringstyle=\color{orange!80!black},
  numbers=left, numberstyle=\tiny\color{gray},
  frame=single, rulecolor=\color{gray!40},
  breaklines=true, showstringspaces=false
}

\begin{document}

\begin{center}
  {\LARGE\bfseries \assignmenttitle}\\[0.4em]
  {\large \coursecode\ \coursename}\\[1em]
  \begin{tabular}{rl}
    Student: & \studentname\ (\studentid) \\
    Lecturer: & \lecturer \\
    Date: & \today
  \end{tabular}
\end{center}
\vspace{1em}

\begin{question}
Solve $x^2 - 5x + 6 = 0$.
\end{question}
\begin{solution}
Factorising gives $(x-2)(x-3)=0$, so $x = 2$ or $x = 3$.
\end{solution}

\begin{question}
A car travels \SI{120}{\kilo\metre} in \SI{1.5}{\hour}. Find its average speed.
\end{question}
\begin{solution}
\[
  v = \frac{d}{t} = \frac{\SI{120}{\kilo\metre}}{\SI{1.5}{\hour}} = \SI{80}{\kilo\metre\per\hour}.
\]
\end{solution}

\begin{question}
Write a Python function that returns the factorial of $n$.
\end{question}
\begin{solution}
\begin{lstlisting}[language=Python]
def factorial(n):
    # Iterative factorial
    result = 1
    for i in range(2, n + 1):
        result *= i
    return result
\end{lstlisting}
\end{solution}

\section*{Lab Results}
\begin{table}[h]
  \centering
  \begin{tabular}{S[table-format=2.1] S[table-format=3.2]}
    \toprule
    {Voltage (\si{\volt})} & {Current (\si{\milli\ampere})} \\
    \midrule
    1.0 & 10.20 \\
    2.0 & 20.35 \\
    5.0 & 50.90 \\
    \bottomrule
  \end{tabular}
  \caption{Measured current for each applied voltage.}
\end{table}

\end{document}
